\documentclass[lualatex, paper={90mm, 160mm}]{jlreq}
\usepackage{tikz}
% luatexja は jlreq[lualatex] が自動的に読み込む
% lltjext: 縦書きminipage (<t>オプション) のために必要
\usepackage{lltjext}
\usetikzlibrary{calc}

\begin{document}
\pagestyle{empty}

\begin{tikzpicture}[remember picture, overlay]

  % ━━━━━━━━━━━━━━━━━━━━━━━━━━━━━━━━━━
  % 1. 背景模様: 放射線 + 同心長方形
  %    中心から放射する細線と、中心を囲む同心長方形を重ねる
  % ━━━━━━━━━━━━━━━━━━━━━━━━━━━━━━━━━━

  % 放射線（10°刻み、36本）
  \begin{scope}[opacity=0.35, color=black!85, very thin]
    \foreach \angle in {0, 10, ..., 350} {
      \draw (current page.center) -- ++(\angle:130mm);
    }
  \end{scope}

  % 同心長方形（ページ縦横比に合わせた長方形を等間隔に重ねる）
  \begin{scope}[opacity=0.40, color=black!90, very thin]
    \foreach \s in {1,...,9} {
      \draw ($(current page.center) - (\s*4.5mm, \s*8mm)$)
            rectangle
            ($(current page.center) + (\s*4.5mm, \s*8mm)$);
    }
  \end{scope}

  % ━━━━━━━━━━━━━━━━━━━━━━━━━━━━━━━━━━
  % 2. 二重枠線
  % ━━━━━━━━━━━━━━━━━━━━━━━━━━━━━━━━━━
  \draw[line width=1.2pt, color=black!80]
    ($(current page.south west) + (6mm, 6mm)$)
    rectangle ($(current page.north east) - (6mm, 6mm)$);
  \draw[line width=0.3pt, color=black!80]
    ($(current page.south west) + (7.5mm, 7.5mm)$)
    rectangle ($(current page.north east) - (7.5mm, 7.5mm)$);

  % ━━━━━━━━━━━━━━━━━━━━━━━━━━━━━━━━━━
  % 3. タイトル（縦書き）
  % ━━━━━━━━━━━━━━━━━━━━━━━━━━━━━━━━━━
  \node[anchor=north] at ($(current page.center) + (15mm, 50mm)$) {
    \begin{minipage}<t>{15em}
      \begin{center}
        {\fontsize{28pt}{40pt}\selectfont 走れメロス\par}
      \end{center}
    \end{minipage}
  };

  % ━━━━━━━━━━━━━━━━━━━━━━━━━━━━━━━━━━
  % 4. 作者名（縦書き）
  % ━━━━━━━━━━━━━━━━━━━━━━━━━━━━━━━━━━
  \node[anchor=north] at ($(current page.center) + (-15mm, -10mm)$) {
    \begin{minipage}<t>{10em}
      \begin{center}
        {\fontsize{16pt}{24pt}\selectfont 太宰 治\par}
      \end{center}
    \end{minipage}
  };

\end{tikzpicture}
\end{document}